\documentclass[aspectratio=169]{beamer}
\usepackage[english]{babel}
\usepackage[utf8]{inputenc}
\usepackage[T1]{fontenc}
\usepackage{lmodern}
\usepackage{listings}
\usepackage{xcolor}
\usepackage{graphicx}
\usepackage{textcomp}
\DeclareUnicodeCharacter{2014}{---}
\DeclareUnicodeCharacter{2013}{--}
\DeclareUnicodeCharacter{2212}{-}
\DeclareUnicodeCharacter{2192}{\ensuremath{\rightarrow}}
\DeclareUnicodeCharacter{00B1}{\ensuremath{\pm}}
\DeclareUnicodeCharacter{00D7}{\ensuremath{\times}}
\DeclareUnicodeCharacter{00B7}{\ensuremath{\cdot}}
\DeclareUnicodeCharacter{20AC}{EUR}
\DeclareUnicodeCharacter{2070}{\textsuperscript{0}}
\DeclareUnicodeCharacter{00B9}{\textsuperscript{1}}
\DeclareUnicodeCharacter{00B2}{\textsuperscript{2}}
\DeclareUnicodeCharacter{00B3}{\textsuperscript{3}}
\DeclareUnicodeCharacter{2074}{\textsuperscript{4}}
\DeclareUnicodeCharacter{2075}{\textsuperscript{5}}
\DeclareUnicodeCharacter{2076}{\textsuperscript{6}}
\DeclareUnicodeCharacter{2077}{\textsuperscript{7}}
\DeclareUnicodeCharacter{2078}{\textsuperscript{8}}
\DeclareUnicodeCharacter{2079}{\textsuperscript{9}}
\DeclareUnicodeCharacter{2248}{\ensuremath{\approx}}
\DeclareUnicodeCharacter{2264}{\ensuremath{\le}}

% Colors for code
\definecolor{codebg}{RGB}{245,245,245}
\definecolor{codecomment}{RGB}{0,128,0}
\definecolor{codekeyword}{RGB}{0,0,255}
\definecolor{codestring}{RGB}{163,21,21}
\definecolor{bandcolor}{RGB}{200,50,50}

\lstdefinestyle{pythonstyle}{
    language=Python,
    backgroundcolor=\color{codebg},
    basicstyle=\ttfamily\scriptsize,
    keywordstyle=\color{codekeyword},
    commentstyle=\color{codecomment},
    stringstyle=\color{codestring},
    showstringspaces=false,
    breaklines=true,
    frame=single,
    rulecolor=\color{black},
    escapeinside={(*@}{@*)},
    moredelim=**[l][\color{bandcolor}\bfseries]{\#\ ---\ NEW},
    moredelim=**[l][\color{bandcolor}\bfseries]{\#\ ---\ CHANGED},
    literate=
      {á}{{\'a}}1 {é}{{\'e}}1 {í}{{\'i}}1 {ó}{{\'o}}1 {ú}{{\'u}}1
      {Á}{{\'A}}1 {É}{{\'E}}1 {Í}{{\'I}}1 {Ó}{{\'O}}1 {Ú}{{\'U}}1
      {ñ}{{\~n}}1 {Ñ}{{\~N}}1 {ü}{{\"u}}1 {Ü}{{\"U}}1
      {¿}{{?`}}1 {¡}{{!`}}1
      {—}{{---}}1 {–}{{--}}1 {−}{{-}}1
      {±}{{\ensuremath{\pm}}}1 {×}{{\ensuremath{\times}}}1
      {→}{{\ensuremath{\rightarrow}}}1
      {°}{{\textdegree}}1
      {·}{{\ensuremath{\cdot}}}1
      {€}{{\texteuro}}1
      {≈}{{\ensuremath{\approx}}}1
      {≤}{{\ensuremath{\le}}}1
      {⁰}{{\textsuperscript{0}}}1 {¹}{{\textsuperscript{1}}}1
      {²}{{\textsuperscript{2}}}1 {³}{{\textsuperscript{3}}}1
      {⁴}{{\textsuperscript{4}}}1 {⁵}{{\textsuperscript{5}}}1
      {⁶}{{\textsuperscript{6}}}1 {⁷}{{\textsuperscript{7}}}1
      {⁸}{{\textsuperscript{8}}}1 {⁹}{{\textsuperscript{9}}}1
}

\lstset{style=pythonstyle}

% Title slide macro: \lessontitle{Lesson Number}{Title}
\newcommand{\lessontitle}[2]{
    \title{Lesson #1: #2}
    \author{}
    \institute{Tec de Monterrey}
    \begin{frame}[plain]
        \titlepage
    \end{frame}
}

% Figure frame macro: \figureframe{Title}{Image path}
\newcommand{\figureframe}[2]{
    \begin{frame}{#1}
        \begin{center}
            \includegraphics[height=0.8\textheight,width=\textwidth,keepaspectratio]{#2}
        \end{center}
    \end{frame}
}


% Compact classroom route. The seven scripts remain the authority.
% Every listing below points to a real band in the canonical source code.
\lstset{
    basicstyle=\ttfamily\fontsize{5.4pt}{6.4pt}\selectfont,
    breaklines=true,
    breakatwhitespace=true,
    columns=fullflexible,
    keepspaces=true,
    xleftmargin=0.2em,
    framesep=2pt,
    aboveskip=0pt,
    belowskip=0pt
}

\newcommand{\resultframe}[3]{%
    \begin{frame}{#1}
        \begin{center}
            \includegraphics[height=0.63\textheight,width=\textwidth,keepaspectratio]{#2}
        \end{center}
        \vspace{0.2em}
        {\small #3\par}
    \end{frame}%
}

\date{}

\begin{document}

\lessontitle{01}{Your First Genetic Algorithm: compact route}

\begin{frame}{The 50-minute route}
    \begin{enumerate}
        \item Formulate the landscape and inspect the best sampled point.
        \item Build a population: candidates before a single answer.
        \item Separate selection, crossover, and mutation.
        \item Combine them for ten generations and change only the seed.
    \end{enumerate}

    \vspace{0.5em}
    \textbf{Guiding question:} What does each operator add that the previous
    operator could not do?
\end{frame}

\begin{frame}{Problem and reference}
    \[
        \max_{x\in[-10,10]} f(x),\qquad f(x)=\sin x-0.2|x|
    \]
    A 400-point grid provides a numerical reference. The genetic algorithm
    sees candidates, evaluates them, and transforms a finite population.

    \vspace{0.6em}
    \textbf{Keep separate:} the best sampled point is not a proof of global
    optimality.
\end{frame}

\resultframe{1. Landscape}{../figures/first_example_01_the_landscape.png}{%
    The landscape has several peaks. A local search can stop on the peak where
    it starts; the grid reports \texttt{x = +1.378} as this demonstration's
    numerical reference.}

\begin{frame}[fragile]{2. Population: preserve alternatives}
    The recipe adds \texttt{Individual}, \texttt{create\_random()}, and the
    initial population. The algorithm still does not select or modify
    candidates: it stores ten evaluated alternatives.

    \vspace{0.15em}
    \lstinputlisting[basicstyle=\ttfamily\fontsize{4.7pt}{5.6pt}\selectfont,
        firstline=58,lastline=59]{../src/first_example_02_random_population.py}
    \lstinputlisting[basicstyle=\ttfamily\fontsize{4.7pt}{5.6pt}\selectfont,
        firstline=71,lastline=82]{../src/first_example_02_random_population.py}
    \lstinputlisting[basicstyle=\ttfamily\fontsize{4.7pt}{5.6pt}\selectfont,
        firstline=109,lastline=110]{../src/first_example_02_random_population.py}
    \lstinputlisting[basicstyle=\ttfamily\fontsize{4.7pt}{5.6pt}\selectfont,
        firstline=129,lastline=141]{../src/first_example_02_random_population.py}
\end{frame}

\resultframe{2. Initial population}{../figures/first_example_02_random_population.png}{%
    With seed 52, the best initial candidate is \texttt{x = -0.323}. Ten blind
    guesses cover more possibilities than one, but they are still guesses.}

\begin{frame}[fragile]{3. Selection: change multiplicities}
    Tournament selection keeps the best of three sampled candidates and repeats
    until the population is full. It creates no new gene: selection only copies.

    \vspace{0.25em}
    \lstinputlisting[basicstyle=\ttfamily\fontsize{5.0pt}{6.0pt}\selectfont,
        firstline=121,lastline=122]{../src/first_example_03_selection.py}
    \lstinputlisting[basicstyle=\ttfamily\fontsize{5.0pt}{6.0pt}\selectfont,
        firstline=142,lastline=151]{../src/first_example_03_selection.py}
\end{frame}

\resultframe{3. Selection}{../figures/first_example_03_selection.png}{%
    The run produces \texttt{4/10} individuals with zero copies. Selection
    concentrates the population, but it cannot explore by itself.}

\begin{frame}[fragile]{4. Crossover: create new values}
    \texttt{clamp()} preserves the domain; \texttt{crossover\_blend()} combines
    two genes; \texttt{crossover()} returns new individuals. The blend
    parameter can leave the parental interval while staying in \([-10,10]\).

    \vspace{0.1em}
    \lstinputlisting[basicstyle=\ttfamily\fontsize{4.5pt}{5.4pt}\selectfont,
        firstline=61,lastline=61]{../src/first_example_04_crossover.py}
    \lstinputlisting[basicstyle=\ttfamily\fontsize{4.5pt}{5.4pt}\selectfont,
        firstline=80,lastline=82]{../src/first_example_04_crossover.py}
    \lstinputlisting[basicstyle=\ttfamily\fontsize{4.5pt}{5.4pt}\selectfont,
        firstline=131,lastline=145]{../src/first_example_04_crossover.py}
    \lstinputlisting[basicstyle=\ttfamily\fontsize{4.5pt}{5.4pt}\selectfont,
        firstline=167,lastline=168]{../src/first_example_04_crossover.py}
    \lstinputlisting[basicstyle=\ttfamily\fontsize{4.5pt}{5.4pt}\selectfont,
        firstline=181,lastline=184]{../src/first_example_04_crossover.py}
\end{frame}

\resultframe{4. Crossover}{../figures/first_example_04_crossover.png}{%
    With parents \texttt{-2} and \texttt{+3}, \texttt{14/20} children fall
    outside the parental interval. New values do not guarantee better fitness.}

\begin{frame}[fragile]{5. Mutation: extend the reach}
    Gaussian mutation adds noise and reapplies the bounds. The parameter
    \texttt{sigma} controls the typical jump size.

    \vspace{0.2em}
    \lstinputlisting[basicstyle=\ttfamily\fontsize{5.0pt}{6.0pt}\selectfont,
        firstline=123,lastline=124]{../src/first_example_05_mutation.py}
    \lstinputlisting[basicstyle=\ttfamily\fontsize{5.0pt}{6.0pt}\selectfont,
        firstline=143,lastline=150]{../src/first_example_05_mutation.py}
    \lstinputlisting[basicstyle=\ttfamily\fontsize{5.0pt}{6.0pt}\selectfont,
        firstline=164,lastline=167]{../src/first_example_05_mutation.py}
\end{frame}

\resultframe{5. Mutation}{../figures/first_example_05_mutation.png}{%
    From a local hill, \texttt{sigma=1.0} produces \texttt{0/2000} arrivals and
    \texttt{sigma=3.0} produces \texttt{110/2000}. Zero observations in a finite
    sample do not prove a probability of zero.}

\begin{frame}[fragile]{6. Generational loop}
    The complete run repeats
    \[
        \text{initialise}\to[\text{select}\to\text{crossover}\to\text{mutate}\to\text{replace}]\times10.
    \]
    The following band assembles the operators into the same reproducible run.

    \vspace{0.2em}
    \lstinputlisting[basicstyle=\ttfamily\fontsize{4.8pt}{5.7pt}\selectfont,
        firstline=200,lastline=218]{../src/first_example_06_the_full_loop.py}
\end{frame}

\resultframe{6. Complete loop}{../figures/first_example_06_the_full_loop.png}{%
    With seed 52, ten generations finish at \texttt{x = +1.372, f = +0.706},
    near the grid reference \texttt{x = +1.378}.}

\begin{frame}[fragile]{7. Change only the seed}
    Everything stays fixed except the random sequence. The change tests run
    sensitivity, not implementation correctness.

    \vspace{0.55em}
    \lstinputlisting[basicstyle=\ttfamily\fontsize{7.2pt}{8.4pt}\selectfont,
        firstline=36,lastline=38]{../src/first_example_07_local_optimum.py}
\end{frame}

\resultframe{7. Local optimum}{../figures/first_example_07_local_optimum.png}{%
    With seed 16, the same algorithm finishes at \texttt{x = -4.417, f = +0.073}.
    A successful run does not prove global optimality.}

\begin{frame}{Conceptual close}
    \begin{itemize}
        \item The population preserves alternatives.
        \item Selection amplifies what already works.
        \item Crossover recombines; mutation restores reach.
        \item The loop produces a trajectory, not a guarantee.
    \end{itemize}

    \vspace{0.5em}
    \textbf{Next:} Lesson 02 names and separates the algorithm's phases.
\end{frame}

\end{document}
